% RESEARCH_DRAFT_SHA256: 0744451656b2eb45e0b91e48406238076dda612ac981906ac46b085e6c02edf8
% Layout and source-reference wrappers expanded for the web renderer.
\documentclass[11pt]{article}
\usepackage[margin=25mm]{geometry}
\usepackage{amsthm}
\usepackage{fontspec,unicode-math}
\setmainfont{Latin Modern Roman}
\setmathfont{Latin Modern Math}
\usepackage{amsmath,amssymb,mathtools}
\usepackage{xurl,hyperref}
\hypersetup{colorlinks=true,urlcolor=blue,linkcolor=blue}
\setcounter{secnumdepth}{0}
\setlength{\parindent}{0pt}
\setlength{\parskip}{5pt}
\setlength{\emergencystretch}{3em}
\newcommand{\sref}[2]{\hyperlink{source:#1:#2}{[S#2]}}
\newcommand{\cataloguescope}[1]{\paragraph{Source question.}#1}
\newcommand{\floor}[1]{\left\lfloor#1\right\rfloor}
\newtheorem{proposition}{Proposition}
\newtheorem{lemma}[proposition]{Lemma}
\begin{document}
\paragraph{Research manuscript.} OpenAI \texttt{gpt-6-astra}, reasoning \texttt{ultra}; 11 October 2026.
\paragraph{Preserved source body.} The complete catalogue mathematical text and sources follow with layout and source links expanded, including source qualifications and its historical preparation notes. New research begins at the section headed Research attempt.

\section{Arbitrary logistic cell gaps with positive margins and ordinary-margin obstructions}\label{statement}

\subsection{Definitions and mathematical statement}

\paragraph{Shared notation.}Standard number systems are denoted by $\mathbb N,\mathbb Z,\mathbb Q,\mathbb R,\mathbb C$; index ranges, computational models and specialized notation follow the cited problem source and the definitions below.


For a hierarchical table model with margins $b$, a cell's attainable set is $S_i(b)=\{x_i:x\in\mathbb Z_{\ge0}^d,Ax=b\}$. A gap is a run of integers between attainable values that are not in this set. Arbitrarily large gaps means that for every positive integer $N$ one can choose allowed margins and a cell with a gap of length at least $N$.
\paragraph{Statement recorded in the supplied source.}
Problem 55 (Karr). Find a set of strictly positive (or strictly > 5) margins for which the set of tables which has these margins has an arbitrarily large gap in a cell entry. Dinwoodie: Are there examples of this phenomenon in logistic regression? \href{https://aimath.org/WWN/compalgstat/compalgstat.pdf}{[S2]}
\paragraph{Scope qualification.}
Both the strictly-positive and the stronger greater-than-five margin variants, and the accompanying logistic-regression question, are retained.
\paragraph{Source status.}
The supplied catalogue labels this record partially solved. Its source-status metadata was checked against the supplied snapshot on 2026-10-11. This import does not independently establish current open status or mathematical correctness.
\subsection{Short English statement}
Can strictly positive table margins force arbitrarily long gaps in the possible values of a cell?
\subsection{Sources}
\begin{itemize}
\item[S1] Ulam AI and the UnsolvedMath Contributors, UnsolvedMath: A Curated Collection of Open Mathematics Problems, catalogue entry. Catalogue curation and metadata: CC BY 4.0. The supplied source question is retained; definitions and the English summary are editorial additions.
\url{https://huggingface.co/datasets/ulamai/UnsolvedMath}.
\item[S2] AIM Workshop Problem Lists, original problem source linked by the catalogue
\url{https://aimath.org/WWN/compalgstat/compalgstat.pdf}.
\end{itemize}
\paragraph{Preparation and review.}
Original question attributed to the cited source and UnsolvedMath contributors. Codex (GPT-6; exact serving revision not exposed) assisted with definitions, English summaries, source inspection and formatting. No new human mathematical authorship or independent proof verification is asserted. Catalogue identity checked against the complete supplied registry and existing repository statements; known flagged repetitions and obvious identical conjectures excluded. Full mathematical-equivalence/literature review has not been performed. Statement preparation and final review followed contribution prompts C and D; automated checks verify structure and identity, not open status.
\label{end:statement}


\section{Research attempt}
\subsection{Literature and scope}
Problem 55 asks for large gaps with positive margins and separately asks about logistic regression. I constructed an exact conditional logistic fiber with odd covariate spacing. It realizes the logistic phenomenon but does not convert the weighted sufficient statistic into ordinary hierarchical margins; the main margin question is not silently replaced.
\begin{proposition}
Fix integers $L\ge1$ and $q\ge6$. For three binomial groups with covariates $(0,2,2L+1)$ and sizes
\[
(n_1,n_2,n_3)=(q(2L-1),q(2L+1),2q),
\]
condition on total successes $q(2L+1)$ and weighted successes $2q(2L+1)$. The feasible success triples are exactly
\[
(s_1,s_2,s_3)=((2L-1)t,(2L+1)(q-t),2t),\qquad t=0,\ldots,q.
\]
Thus the second-cell spectrum has gaps of $2L$ missing integers between consecutive attainable values, although every row and column margin of the associated success/failure table exceeds five.
\end{proposition}
\begin{proof}
The weighted equation $2s_2+(2L+1)s_3=2q(2L+1)$ forces $s_3$ even. Write $s_3=2t$ and solve it for $s_2=(2L+1)(q-t)$. The total-success equation then gives $s_1=(2L-1)t$. The bounds $0\le s_i\le n_i$ are equivalent to $0\le t\le q$, so all and only the listed triples occur. Consecutive attainable $s_2$ values differ by $2L+1$, leaving $2L$ integers between them.

Column margins are the group sizes and are at least six. Their sum is $q(4L+2)=2q(2L+1)$, so both success and failure row totals equal $q(2L+1)>5$. In logistic regression with log odds $\alpha+\beta x_i$, the joint mass at $s$ is proportional, apart from fixed binomial denominators, to $\prod_i\binom{n_i}{s_i}\exp(\alpha\sum_i s_i+\beta\sum_i x_is_i)$. Conditioning on the two displayed sums removes $\alpha,\beta$ and gives exactly this fiber.
\end{proof}
\subsection{Verification and remaining gaps}
For $L=4,q=6$, exhaustive bounded enumeration found exactly seven triples, with second counts $0,9,18,27,36,45,54$. The gap is eight and group sizes are $42,54,12$. The mechanism is an elementary lattice congruence, not a new unrestricted gap theorem. Without the extra weighted statistic, these two-way margins do not force this spectrum. That is the precise obstruction to treating the construction as a solution for ordinary hierarchical margins in the first part of the source problem.

\section{References and source attribution}
The original question is credited to the AIM workshop list prepared by Seth Sullivant (2004), \url{https://aimath.org/WWN/compalgstat/compalgstat.pdf}. The source question and catalogue attribution are reproduced above. This list was accessed on 11 October 2026; it is historical evidence of the question, not a certification of present open status.


\paragraph{Integration review.} Codex (GPT-6; exact serving revision not exposed) reviewed the written arguments and expanded layout and source-reference wrappers. Repository metadata and the production web renderer were checked using Python 3.13 and JavaScript with MathJax 3.2.2. This model-assisted review is not independent human proof verification; no PDF compilation or formal proof checking is asserted.

\section{Completion Estimate}
\noindent\textbf{COMPLETION: 0\%}
No substantiated new progress toward the entire original target is claimed. The worked results and failed route are delimited above; a known result, illustrative model, or auxiliary obstruction is not counted as original-target completion.

\section{Final status and attribution}
\textbf{Attempt status: unresolved.} This describes this research attempt, not a new ruling on the source catalogue's status or on external literature claims.
The mathematics and this self-audit were generated by OpenAI \texttt{gpt-6-astra}, reasoning \texttt{ultra}, on 11 October 2026. Codex (GPT-6; exact serving revision not exposed) orchestrated the batch and contributed formatting/repository review. No human mathematical authorship, independent human proof checking, or formal verification is claimed. The historical preparation paragraphs in the reproduced catalogue body concern that body only.

Contribution prompt A was applied to the complete definition and supplied source record, followed by the adversarial checks of prompt B. Local reproduction files are \texttt{batch005/checks.py}, \texttt{computations.json}, and the per-ID audit JSON. They were not public links at preparation time. Numerical checks are finite and do not prove unrestricted results. No LaTeX engine was found locally; compilation was not performed by this research agent.

\end{document}
